\section{总结与延伸}

这堂课最值得保留的不是某个 speedup headline，而是一套把 architecture、statistics 与 systems 连起来的思考方式：如果能从数据统计得到可解释初值，就可能减少随机训练；如果底层表示固定，就可能缓存重复比较；如果任务极窄且数据本地生成，就可能把训练移到 edge。每一步都带来新约束，不能只取收益、丢掉前提。

\subsection{三层压缩：初始化、计算复用与本地数据闭环}

第一层是 initialization：bit-cipher 提供 deterministic non-negative representation，generalized co-occurrence 为 softmax layer 提供 warm start。第二层是 computation reuse：frozen embeddings 使 pairwise comparisons 可缓存，dynamic context 避免短样本支付最长窗口成本。第三层是 data loop：edge controller 从真实交互中采集 transcript/action pairs，用窄任务换取少数据可学性。

\begin{importantbox}{课程的最小可信结论}
\begin{enumerate}
  \item 对满足假设的 softmax layer，data-derived warm start 可以优于随机起点。
  \item Warm start 后的 iterative training 仍有价值，尤其对多层 composition 和最终表示适配。
  \item Fixed representations 让部分比较可缓存，但 cache、内存与 invalidation 必须计入系统成本。
  \item Tiny local controller 可以从应用交互学习窄 action，但不能据此推出通用对话能力。
\end{enumerate}
\end{importantbox}

\subsection{复现实验 checklist}

复现这类工作时，最容易遗漏的是假设与额外成本。下面的 checklist 应和结果图一起发布。

\begin{lstlisting}[caption={SAFFU/PLM 复现与发布清单}]
record_official_architecture_and_parameter_shapes()
verify_nonnegative_features_and_softmax_activation()
publish_explicit_initialization_formula_and_epsilon()
report_priming_number_distribution_and_sensitivity()
separate_stats_cache_backprop_and_evaluation_time()
compare_tuned_cold_warm_frozen_and_thawed_baselines()
document_cache_size_location_hit_rate_and_invalidation()
evaluate_heldout_users_paraphrases_noise_and_class_balance()
report_end_to_end_edge_latency_not_model_latency_only()
publish_code_data_splits_seeds_and_failure_cases()
\end{lstlisting}

\begin{warningbox}{看到“无需 pretraining/backprop”时立即追问}
是哪一层无需？显式统计量用了多少数据？后续是否仍做 gradient tuning？Embedding 是否冻结？测试是 training fit 还是 held-out generalization？Unsupported layers 是否随机初始化？若这些问题没有答案，headline 远大于证据。
\end{warningbox}

\subsection{自测问题}

\begin{multicols}{2}
\small
\begin{enumerate}
  \item 标准 $QK^\top$ attention 与 $XWX^\top$ 在参数化上有什么关系？
  \item 为什么 SAFFU 仍然是有参数 attention，而不是无参数相似度？
  \item Near-shallow 描述的是 depth，还是在线更新/计算依赖？
  \item Bit-cipher 为什么需要 $b\ge\lceil\log_2(V+1)\rceil$？
  \item Generalized co-occurrence $H^\top Y$ 与传统词窗共现有什么不同？
  \item 显式 softmax solution 为什么需要 non-negative features？
  \item Priming number $K$ 在公式中修正什么？
  \item 为什么 attention 的 hidden targets 比 decoder targets 更难获得？
  \item Warm-start PPL 图为什么不能直接证明 SAFFU attention 更好？
  \item MNIST 结果能证明什么，又不能证明什么？
  \item Frozen embedding 为什么同时影响训练速度和 cache correctness？
  \item Packing 与 dynamic context 分别减少哪种浪费？
  \item PLM 的 precision 为什么不是 FP16/BF16？
  \item Warm--freeze--thaw 三阶段各自优化什么？
  \item Le Potato 的端到端 latency 包含哪些组件？
  \item Binary switch accuracy 为什么可能被 class imbalance 误导？
  \item 哪些结论在课堂结束时仍只是 future work？
\end{enumerate}
\end{multicols}

\subsection{拓展阅读}

\begingroup
\small
\setlength{\itemsep}{2pt}
\begin{itemize}
  \item Williams and Zhao, \href{https://arxiv.org/abs/2311.07510}{\emph{Explicit Foundation Model Optimization with Self-Attentive Feed-Forward Neural Units}}：SAFFU 架构、显式初始化与 ablation 主论文。
  \item Williams and Zhao, \href{https://arxiv.org/abs/2311.07498}{\emph{Reducing the Need for Backpropagation and Discovering Better Optima With Explicit Optimizations of Neural Networks}}：单层 softmax 显式解、MNIST 与 local multi-layer warm start。
  \item Zhao and Williams, \href{https://arxiv.org/abs/2311.11012}{\emph{Bit Cipher}}：deterministic bit vectors、co-occurrence aggregation 与 transfer experiments。
  \item Williams and Heidenreich, \href{https://arxiv.org/abs/2205.00148}{\emph{To Know by the Company Words Keep and What Else Lies in the Vicinity}}：Word2Vec softmax、co-occurrence factorization 与 radial context 动机。
  \item Frankle and Carbin, \href{https://arxiv.org/abs/1803.03635}{\emph{The Lottery Ticket Hypothesis}}：理解本讲“减少随机抽奖”动机的原始背景，但不要把它当作 SAFFU 的实验结论。
  \item Warstadt et al., \href{https://arxiv.org/abs/2310.00915}{\emph{Findings of the BabyLM Challenge}}：检查 developmentally plausible data scale 与标准 evaluation 的具体要求。
\end{itemize}
\endgroup

\end{document}
